\begin{mistake}{2026-06-16}{01}

\begin{problemcontent}
设函数由方程 \(2y^3 - 2y^2 + 2xy - x^2 = 1\) 所确定，试求 \(y = y(x)\) 的驻点，并判别它是否为极值点。
\end{problemcontent}

\begin{solutioncontent}
方程两边对 \(x\) 求导：
\[ (6y^2 - 4y + 2x)y' + 2y - 2x = 0 \]
令一阶导数 \(y' = 0\)，得 \(y = x\)。

将 \(y = x\) 代回原方程，得：
\[ 2x^3 - x^2 - 1 = 0 \]
因式分解得 \((x - 1)(2x^2 + x + 1) = 0\)，唯一实数根为 \(x = 1\)，此时 \(y = 1\)。故驻点为 \(x = 1\)（对应点 \((1,1)\)）。

将一阶导数等式两边再次对 \(x\) 求导：
\[ \frac{d}{dx}(6y^2 - 4y + 2x) \cdot y' + (6y^2 - 4y + 2x)y'' + 2y' - 2 = 0 \]
代入 \(x=1, y=1, y'=0\)：
\[ 4y'' - 2 = 0 \implies y''(1) = \frac{1}{2} \]
因为驻点处 \(y'(1) = 0\) 且二阶导数 \(y''(1) > 0\)，所以 \(x = 1\) 是该函数的极小值点。
\end{solutioncontent}

\mistakeknowledge{
\begin{itemize}
  \item \textbf{隐函数求二阶导技巧}：直接对一阶导数等式两端再次求导，并立刻代入驻点条件 \(y'=0\)，可以消去大量含 \(y'\) 的项，避免解出 \(y'\) 再用商的求导法则带来的繁琐计算。
  \item \textbf{极值第二充分条件}：若 \(y'(x_0) = 0\)，当 \(y''(x_0) > 0\) 时为极小值点，当 \(y''(x_0) < 0\) 时为极大值点。
  \item \textbf{高次方程求根}：三次及以上多项式方程通常有整数根（如 \(0, \pm1, \pm2\)），试根后用因式分解降次即可。
\end{itemize}}

\end{mistake}
